\centerline{\Huge \bf  Подобие}
\centerline{\Huge \bf Теорема Фалеса}

\begin{flushright}
    \scriptsize{Кира Йошикаге использовал теорему Фалеса посреди боя,\\
                дабы вычислить расстояние до цели.}
\end{flushright}

\problem{1} На сторонах $AB$ и $BC$ треугольника $ABC$ взять точки $P$ и $Q$ таким образом, что $PQ \parallel AC$. Докажите, что медиана $BM$ делит отрезок $PQ$ пополам. 

\problem{2} На сторонах $AB$ и $BC$ треугольника $ABC$ взять точки $P$ и $Q$ таким образом, что $PQ \parallel AC$ и $PQ = \dfrac{1}{2}AC$. Докажите, что любая чевиана$\footnote{Чевиана — отрезок в треугольнике, соединяющий вершину треугольника с точкой на противоположной стороне}$ $BM$ делится отрезком $PQ$ пополам.

\problem{3} В треугольнике ABC проведены высоты $AA'$ и $BB'$ Докажите, что треугольник $A'B'C$ подобен треугольнику $ABC$ и $CB' \cdot CA = CA' \cdot CB.$

\problem{4} Точка $M -$ середина стороны AB паралеллограмма $ABCD$. Найдите, в каком отношении отрезок $CM$ делит диагональ $BD$.

\problem{5} \textit{Золотое правило трапеции.} Середины оснований, точка пересечения диагоналей трапеции и точка пересечения продолжений боковых сторон лежат на одной прямой.

\problem{6} Постройте в трапеции отрезок, параллельный основаниям, который делится ее диагоналями на три равные части.

\problem{7} В угле отмеченна произвольная точка. Проведите через эту точку такую прямую, что отрезок, проходящий через эту точку и ограниченный сторонами угла, делится нашей точкой в заданном отношении.

\textbf{Теорема Фалеса.} Если на одной из двух прямых отложить последовательно несколько равных между собой отрезков и через их концы провести параллельные прямые, пересекающие вторую прямую, то они отсекут на второй прямой равные между собой отрезки.

Более общая формулировка, также называемая теоремой о пропорциональных отрезках
Параллельные секущие образуют на прямых пропорциональные отрезки:

\hspace{3mm}

${\displaystyle {\frac {A_{1}A_{2}}{B_{1}B_{2}}}={\frac {A_{2}A_{3}}{B_{2}B_{3}}}={\frac {A_{1}A_{3}}{B_{1}B_{3}}}.}$

\hspace{3mm}

\textbf{Обратная теорема Фалеса.}
Если прямые, пересекающие две другие прямые (параллельные или нет), отсекают на обеих из них равные (или пропорциональные) между собой отрезки, начиная от вершины, то такие прямые параллельны.